% CSC 6740 class presentation (Beamer), 10 x 5.625 in, laid out like the Google Slides version.
% Story: A walls move -> B the motion is stochastic -> C OU model -> D data -> E fitted parameters ->
% F a simple stochastic simulation. Model evaluation is in the appendix (numbered separately).
% Build (from this folder): latexmk -pdf deck.tex   (speaker notes on a second screen: latexmk -pdf deck_notes.tex)
% Figures: figures/ (finished images; see README.md for their provenance).
% The current version of the talk, with animations and the travel-time results, is class_deck_2026-11-18.pptx.
\documentclass[t]{beamer}
\geometry{papersize={10in,5.625in}}
\ifdefined\shownotes
  \setbeameroption{show notes on second screen=right}
\fi
\usepackage[T1]{fontenc}
\usepackage{helvet}                      % Helvetica: the metric twin of Arial
\renewcommand{\familydefault}{\sfdefault}
\usepackage{amsmath,amssymb}
\usepackage{array}
\usepackage{appendixnumberbeamer}
\usepackage[absolute,overlay]{textpos}
\setlength{\TPHorizModule}{1in}
\setlength{\TPVertModule}{1in}
\textblockorigin{0in}{0in}
\graphicspath{{figures/}}

% ---- Google deck look -----------------------------------------------------------------------------
\definecolor{bul}{HTML}{595959}
\definecolor{cap}{HTML}{595959}
\setbeamertemplate{navigation symbols}{}
\setbeamercolor{normal text}{fg=black,bg=white}
\setbeamertemplate{frametitle}{}
\setbeamertemplate{footline}{\hfill{\color{cap}\fontsize{8}{9}\selectfont
  \insertframenumber\,/\,\inserttotalframenumber}\hspace*{0.25in}\vspace*{0.12in}}
\newcommand{\gsize}[1]{\fontsize{#1}{\the\dimexpr #1pt*115/100\relax}\selectfont}
% slide titles: full-width at (0.43, 0.30) in; left-column at (0.43, 0.72) in, 3.05 in wide
\newcommand{\gtitle}[1]{\begin{textblock}{9.1}(0.43,0.36)\fontsize{20}{24}\selectfont\textcolor{black}{#1}\end{textblock}}
\newcommand{\gtitleL}[1]{\begin{textblock}{3.05}(0.43,0.76)\raggedright\fontsize{20}{24}\selectfont\textcolor{black}{#1}\end{textblock}}
% bullets: gray dot at 7.5 pt, text at 22 pt, 6 pt above each paragraph, 1.15 line spacing
\newenvironment{gitems}[1][15]{\gsize{#1}\begin{list}{\makebox[14.5pt][l]{\textcolor{bul}{\textbullet}}}%
  {\setlength{\leftmargin}{22pt}\setlength{\labelwidth}{14.5pt}\setlength{\labelsep}{0pt}%
   \setlength{\itemsep}{6pt}\setlength{\parsep}{0pt}\setlength{\topsep}{0pt}\setlength{\itemindent}{0pt}}%
  \raggedright}{\end{list}}
% \gbox{x}{y}{w}{content}, \gimg{x}{y}{w}{h}{file}, \gimgw{x}{y}{w}{file}, \geqn{x}{y}{w}{math}, \gcap{x}{y}{w}{text}
\newcommand{\gbox}[4]{\begin{textblock}{#3}(#1,#2)#4\end{textblock}}
\newcommand{\gimg}[5]{\begin{textblock}{#3}(#1,#2)\includegraphics[width=#3in,height=#4in]{#5}\end{textblock}}
\newcommand{\gimgw}[4]{\begin{textblock}{#3}(#1,#2)\includegraphics[width=#3in]{#4}\end{textblock}}
\newcommand{\geqn}[4]{\begin{textblock}{#3}(#1,#2)\centering\fontsize{17}{21}\selectfont$\displaystyle #4$\end{textblock}}
\newcommand{\gcap}[4]{\begin{textblock}{#3}(#1,#2)\centering\fontsize{11}{13}\selectfont\textcolor{cap}{#4}\end{textblock}}
\newcommand{\kT}{k_B T}
\newcolumntype{L}[1]{>{\raggedright\arraybackslash}p{#1}}
\newcolumntype{C}[1]{>{\centering\arraybackslash}p{#1}}

\begin{document}

% =================================================================================================== title
\begin{frame}[plain]
\gbox{0.5}{1.82}{9.0}{\centering\fontsize{30}{36}\selectfont Stochastic Cells\\ Diffusion Inside a Randomly Moving Membrane}
\gbox{0.5}{3.12}{9.0}{\centering\fontsize{16}{19}\selectfont Tommy Stell, Georgia State University}
\note{(20 s) I'm Tommy Stell. This project mines public high-speed videos of red blood cells. The story: cell
walls move; the motion is random; a simple stochastic model describes it; the data we have; the parameters the
data give; and a simple simulation of molecules inside the moving wall.}
\end{frame}

% =================================================================================================== A
\begin{frame}
\gtitle{Cell walls move, but cell models fix them}
\gbox{0.43}{0.95}{9.1}{\begin{gitems}
\item Models of transport and signaling inside a cell solve their equations in a fixed shape
\item Real membranes flicker: about 30\,nm rms in red blood cells, 24\,nm once the cell is held in a trap
      (Caselli 2024); the amplitude rises 1.4$\times$ with ATP (Monzel 2015)
\item Nanometres can matter: touch receptors gate their ion channels at $\sim$13\,nm of displacement (Poole 2014)
\item Of 38 experimental papers we reviewed, none tests whether this motion changes transport inside the cell
\item Public 1\,kHz flicker videos let us describe the wall statistically, fit a model to it, and simulate what
      fixing it costs
\end{gitems}}
\note{(55 s) Almost every model of transport or signaling inside a cell solves its equations in a fixed shape.
But membranes are never still. A red blood cell's edge moves about 30 nanometres; less once it is held in an
optical trap, more when the cell has ATP. That sounds negligible, but touch receptors open their ion channels at
about 13 nanometres of displacement, so nanometre motion is not automatically harmless. Of 38 experimental papers
we reviewed, none tests whether this motion changes transport inside the cell. Either answer is useful: if it
matters, fixed-shape models are biased; if it does not, we can say by how much and where that stops being true.}
\end{frame}

\begin{frame}
\gtitle{Two questions about a randomly moving cell wall}
\gbox{0.43}{0.95}{9.1}{\begin{gitems}
\item Q1, inverse: what can contour recordings tell us about the wall?
\item Q2, forward: how much does wall motion change transport inside the cell?
\item Pipeline: describe the wall statistically, fit a stochastic model, simulate transport inside it
\end{gitems}}
\gimg{0.77}{2.67}{8.46}{2.79}{pipeline_2d.png}
\note{(40 s) Two questions. Inverse: what can the recordings tell us about the wall? Forward: does its motion
matter inside? The figure is the pipeline: track the contour, describe it as a random field, ask which features
of that field are identifiable, then simulate diffusion inside it. The transport results are computed in a
two-dimensional equatorial section.}
\end{frame}

% =================================================================================================== B
\begin{frame}
\gtitleL{The wall moves at random}
\gbox{0.43}{1.6}{3.05}{\begin{gitems}
\item One red blood cell's contour for one second, filmed at 1\,kHz: red outward, blue inward
\item The motion is irregular in angle and in time
\item Raw wall motion is 23--32\,nm rms per contour point, under 1\% of the 4\,$\mu$m radius
\end{gitems}}
\gimg{3.75}{0.81}{5.85}{4.03}{gs_kymograph.png}
\note{(45 s) This is the measured wall of one cell over one second, unrolled around the contour: red where it
bulges out, blue where it moves in. There is no regular pattern in angle or in time. The motion is tens of
nanometres, under one percent of the radius.}
\end{frame}

\begin{frame}
\gtitle{It moves as a Gaussian process with memory}
\gimgw{1.0}{0.97}{8.0}{stochastic_stats.png}
\gbox{0.43}{4.0}{9.1}{\begin{gitems}[14]
\item (a) Displacements are close to Gaussian: kurtosis 3.01--3.39 across the 11 videos (3 for a Gaussian)
\item (b) Shape modes decorrelate over tens of ms, roughly exponentially: the signature of an
      Ornstein--Uhlenbeck (OU) process
\end{gitems}}
\note{(50 s) Two statistics say what kind of random. Left: all displacements in units of their standard
deviation lie close to a Gaussian; kurtosis is 3.0 to 3.4 per video, against 3 for a Gaussian. Right: the
autocorrelation of the largest shape mode decays over tens of milliseconds, close to an exponential in the cell
shown. A Gaussian process with exponentially decaying memory is an Ornstein-Uhlenbeck process, the model on the
next slide. Some cells keep slow correlation past 100 ms; that misfit comes back later.}
\end{frame}

% =================================================================================================== C
\begin{frame}
\gtitleL{Model: each shape mode is an OU process}
\gbox{0.43}{1.6}{3.05}{\begin{gitems}
\item Each shape mode relaxes at a rate set by bending, tension, confinement and viscosity
\item P: thermal noise only
\item A: adds a forcing with memory $\tau_a$ and strength $a$, shared by all modes
\item A is a colored effective temperature; nothing in it shows the forcing is ATP-driven
\end{gitems}}
\geqn{4.0}{0.38}{5.6}{du_{lm}=-r_l\,u_{lm}\,dt+\sqrt{2D_l}\,dW\ \ (+\,f\,dt\ \text{in A})}
\geqn{4.0}{0.88}{5.6}{S_l(\omega)=\frac{2D_l}{\omega^2+r_l^2}\Big[1+\underbrace{\frac{a}{1+\omega^2\tau_a^2}}_{\text{model A only}}\Big]}
\gimg{4.21}{1.86}{4.93}{3.45}{gs_spectrum.png}
\note{(60 s) I write the cell as a sphere plus spherical-harmonic wrinkles u-lm. Each wrinkle is an OU process:
pulled back at rate r-l, set by bending, tension, confinement and viscosity, and kicked by thermal noise. That is
model P. Model A adds a force f that is itself an OU process, with memory tau-a and strength a relative to the
thermal noise; it multiplies the passive Lorentzian by one plus a colored term, which adds power below one over
tau-a, as in the plot. A is really a colored effective temperature.}
\end{frame}

\begin{frame}
\gtitleL{The camera sees one slice of the shape}
\gbox{0.43}{1.6}{3.05}{\begin{gitems}
\item The microscope sees the equatorial contour, not the 3D shape
\item Each contour mode $q$ sums all 3D modes with $l\ge q$ and $l+q$ even
\item So even a passive cell mixes many relaxation rates in each contour mode
\end{gitems}}
\geqn{4.0}{0.42}{5.6}{h_q(t)=R\!\!\sum_{l\ge q,\ l+q\ \text{even}}\!\! Y_{lq}\big(\tfrac{\pi}{2},0\big)\,u_{lq}(t)}
\gimg{3.80}{1.34}{5.74}{3.96}{gs_projection.png}
\note{(40 s) The camera sees only the equator. The bars show how much each 3D mode feeds contour modes 2 and 3;
the weights level off, so every contour mode is a sum of many relaxation rates. I fit the exact observed
spectrum, including 1 ms exposure averaging, tracker smoothing and white noise, with a Whittle likelihood.}
\end{frame}

% =================================================================================================== D
\begin{frame}
\gtitleL{Data: 11 public videos, tracked and validated}
\gbox{0.43}{1.6}{3.05}{\begin{gitems}
\item Caselli et al.\ (2024): 1\,kHz for 5\,s; 6 free, 5 trapped, 4 matched pairs, 7 cells
\item Contour on 2048 rays; synthetic frames: 96\% of the motion recovered, noise far below the signal
\item Noise level calibrated on 25 particles stuck to the glass
\item Features: Fourier contour modes $q=2$--6
\end{gitems}}
\geqn{4.0}{0.42}{5.6}{y(\theta_j,t_k)=\mathcal{S}_\theta\Big[\frac{1}{\Delta}\int_{t_k}^{t_k+\Delta}\mathcal{G}\big(h(\cdot,s)\big)\,ds\Big]+\zeta_{jk}}
\gimg{3.75}{1.28}{5.85}{3.9}{gs_tracker.png}
\gcap{3.75}{5.2}{5.85}{$\mathcal{G}$: image formation; $\Delta$ = 1\,ms exposure; $\mathcal{S}_\theta$: tracker smoothing; $\zeta$: white tracking noise}
\note{(60 s) The data are the public deposit behind Caselli et al. 2024: eleven videos, six free cells and five
held in an optical trap, four of them the same cell both ways. The equation is what a frame measures: the wall,
blurred by the optics, averaged over the 1 ms exposure, smoothed by the tracker, plus white noise. Video has no
ground truth, so I rendered synthetic frames from each video's own intensity profile and noise: the tracker
recovers 96 percent of imposed motion, and its noise is far below the signal in the modes we use. I calibrate
the noise on dust stuck to the glass.}
\end{frame}

% =================================================================================================== E
\begin{frame}
\gtitle{Fitted parameters: the two models side by side}
\gbox{0.43}{0.95}{9.1}{\centering\fontsize{14}{17}\selectfont\renewcommand{\arraystretch}{1.45}%
\begin{tabular}{|L{2.25in}|L{2.95in}|L{3.2in}|}
\hline
 & \textbf{Passive OU (P)} & \textbf{Colored OU (A)} \\ \hline
Bending $\kappa$ & 127--314 $\kT$ \ (11/11) & not identifiable \ (0/11) \\ \hline
Viscosity $\eta$ & 0.008--0.030 Pa\,s \ (11/11) & collapses toward 0 in some cells \\ \hline
Tension $\sigma R^2/\kT$ & not identifiable \ (0/11) & --- \\ \hline
Forcing memory $\tau_a$ & --- & 12--60 ms \ (11/11) \\ \hline
Forcing strength $a$ & --- & 53--3200 \ (11/11) \\ \hline
Data prefer & & A in all 11: $\Delta$AIC $+796$ to $+1479$ \\ \hline
\end{tabular}}
\gcap{0.43}{3.66}{9.1}{Ranges over the 11 videos. ($k$/11): videos whose 95\% profile-likelihood interval closes.
P with temperature fixed at 37\,$^\circ$C. Refitted passive simulations: $\Delta$AIC 95\% quantiles 27--554.}
\gbox{0.43}{4.12}{9.1}{\begin{gitems}[14]
\item P gives clean moduli, but it is the model the data reject. A pins the forcing (memory $\approx$30\,ms,
      $a\gg1$: the slow flicker is mostly not thermal)
\item Under A, absolute stiffness trades off against the noise level, so only ratios are identified; A still
      misfits below 30\,Hz (appendix)
\end{gitems}}
\note{(80 s) Both fits side by side. The passive OU model, with temperature fixed at body temperature, returns a
clean bending modulus, 130 to 310 kT, and an effective viscosity of about 10 to 40 times that of water, with
closed 95 percent intervals in all eleven videos; tension is not identifiable. But the data reject this model:
the colored model wins in all eleven videos by 800 to 1500 AIC units, far beyond what refitted passive
simulations produce. The colored model pins down the forcing: a memory of about 30 milliseconds, 12 to 60, and a
strength much larger than one, so the slow part of the flicker is mostly not thermal. It does not pin down
stiffness: scaling stiffness, viscosity and noise level together leaves the spectra unchanged. And it still
misfits below 30 Hz, so its values are best-fit values of an approximate model.}
\end{frame}

% =================================================================================================== F
\begin{frame}
\gtitleL{Prediction: an effect of about 0.1\%}
\gbox{0.43}{1.6}{3.05}{\begin{gitems}
\item A molecule enters through one channel; $\Phi$ is the share within $0.3R$ of the centre
\item The enclosed cytosol is carried by the wall, so a fluid parcel moves back and forth rather than drifting
\item Estimate: $\Delta\Phi/\Phi\approx20(\sigma/R)^2\approx0.1\%$ for a red blood cell
\item 1\% would take $\sim$90\,nm of wall motion
\end{gitems}}
\geqn{4.0}{1.95}{5.6}{\int_0^t\mathbf{v}\,ds=\sum_k U_k\big[h_k(t)-h_k(0)\big]}
\geqn{4.0}{3.12}{5.6}{\frac{\Delta\Phi}{\Phi}\sim\frac{\mathrm{var}(h)\,|U|^2}{DT}\sim20\Big(\frac{\sigma_{\delta h}}{R}\Big)^2\approx10^{-3}}
\gcap{4.0}{4.2}{5.6}{$T=0.05\,R^2/D$; a heuristic second-order estimate made before any simulation, not a theorem}
\note{(45 s) Now the forward problem. A molecule enters through one channel and I measure Phi, the share near
the centre after a fixed time. The cytosol is enclosed and incompressible, so the wall drags it along and a
parcel just moves back and forth; its displacement is set by the change in the wall, not its history. So the
effect should go through the wall's variance: 20 times sigma over R squared, about 0.1 percent for a red blood
cell, made before any simulation.}
\end{frame}

\begin{frame}
\gtitle{A simple simulation: the fitted wall, with molecules inside}
\gimgw{0.7}{0.95}{8.6}{simple_simulation.png}
\gbox{0.43}{3.72}{9.1}{\begin{gitems}[14]
\item (a, b) Modes $q=2$--6 simulated from model A fitted to the same cell: rms 23\,nm, against 21\,nm measured
\item (c) Molecules enter at a channel, diffuse, and reflect off the moving wall (toy: motion $\times10$, still fluid)
\item Full computation, 11 cells $\times$ 12 walls, cytosol carried: moving vs frozen wall changes $\Phi$ by
      $-0.03\pm0.01\%$
\end{gitems}}
\note{(70 s) Left, the measured wall restricted to the five modes we model; middle, a wall simulated from the
colored OU model fitted to the same cell, on the same colour scale. The rms agrees to 2 nm. The simulation looks
busier in time because the data are averaged over the 1 ms exposure and smoothed by the tracker. Right, the
simulation idea: molecules enter through one channel and diffuse, reflected by the moving wall, here with the
motion exaggerated ten times so you can see it. In the full computation, over 11 cells and 12 simulated walls
each, with the cytosol carried along by the wall, moving versus frozen walls change Phi by three hundredths of a
percent.}
\end{frame}

\begin{frame}
\gtitle{Cytosol carried: wall motion changes transport by about 0.1\% or less}
\gbox{1.7}{1.2}{6.6}{\fontsize{15}{18}\selectfont\renewcommand{\arraystretch}{1.75}%
\begin{tabular}{|L{2.8in}|C{1.6in}|C{1.6in}|}
\hline
\textbf{Change in $\Phi$} & \textbf{$D$ = 5 $\mu$m$^2$/s} & \textbf{$D$ = 0.5 $\mu$m$^2$/s} \\ \hline
Wall motion, cytosol carried & $-0.03\pm0.01$\% & $-0.02\pm0.01$\% \\ \hline
Carried vs still medium & $-0.2\pm0.2$\% & $-0.6\pm0.1$\% \\ \hline
Thermal part of the flicker & $+0.05\pm0.04$\% & $-0.03\pm0.04$\% \\ \hline
Wall motion, still medium & $+0.3\pm0.4$\% & $+1.3\pm0.4$\% \\ \hline
\end{tabular}}
\gbox{1.7}{3.6}{6.6}{\centering\fontsize{13}{15}\selectfont\textcolor{cap}{Pooled over 11 cells, 12 wall paths each,
extrapolated grid; statistical errors}}
\note{(40 s) Pooled over 11 cells and 12 wall paths, after grid extrapolation. With the cytosol carried, wall
motion changes transport by a few hundredths of a percent; carried versus still cytosol by under 1 percent. On
our first grid the medium effect looked like 2 to 3 percent, which was mostly discretization error. Walls taken
straight from the tracked videos give the same answer.}
\end{frame}

% =================================================================================================== end
\begin{frame}
\gtitle{Conclusions}
\gbox{0.43}{1.0}{9.1}{\begin{gitems}[16]
\item Red blood cell walls move by 20--30\,nm, as a Gaussian process with memory of tens of ms
\item An OU model per shape mode describes them, but needs a colored, non-thermal forcing with memory
      $\tau_a\approx12$--60\,ms
\item Contour data identify that forcing and ratios of the moduli, not absolute stiffness
\item Simulated with the fitted model, wall motion changes transport inside the cell by about 0.1\% or less;
      about 90\,nm of motion would be needed for 1\%
\item Next: an ATP-depleted control, and a discocyte rather than a spherical shape
\item Code and data pipeline: github.com/tlstellGSU/brownian\_cells
\end{gitems}}
\note{(45 s) Red blood cell walls move by tens of nanometres as a Gaussian process with memory. An OU model per
shape mode describes that, but it needs a non-thermal colored forcing with a memory of tens of milliseconds. The
contours identify that forcing and ratios of the moduli, not stiffness itself. Simulating molecules inside walls
drawn from the fitted model, the motion changes transport by about 0.1 percent or less: for red blood cells a
fixed wall is a good approximation, and it would take walls moving about 90 nm to reach one percent. Next: an
ATP-depleted control and a realistic shape. Thank you.}
\end{frame}

\begin{frame}
\gtitle{References}
\gbox{0.43}{0.95}{9.1}{\fontsize{10}{11.5}\selectfont\raggedright
\setlength{\parskip}{5pt}
N. Caselli, M. Garc\'ia-Verdugo, M. Calero, et al. Red blood cell flickering activity locally controlled by
holographic optical tweezers. iScience, 2024.\par
N. Caselli. Red blood cells TIF images (data deposit, CC BY 4.0). 2024.\par
C. Monzel, D. Schmidt, C. Kleusch, et al. Measuring fast stochastic displacements of bio-membranes with dynamic
optical displacement spectroscopy. Nature Communications, 2015.\par
K. Poole, R. Herget, L. Lapatsina, H.-D. Ngo, and G. R. Lewin. Tuning Piezo ion channels to detect
molecular-scale movements relevant for fine touch. Nature Communications, 2014.\par
S. T. Milner and S. A. Safran. Dynamical fluctuations of droplet microemulsions and vesicles. Physical Review A,
1987.\par
H. A. Faizi, C. J. Reeves, V. N. Georgiev, P. M. Vlahovska, and R. Dimova. Fluctuation spectroscopy of giant
unilamellar vesicles using confocal and phase contrast microscopy. Soft Matter, 2020.\par
P. Whittle. Estimation and information in stationary time series. Arkiv f\"or Matematik, 1953.\par
H. Akaike. A new look at the statistical model identification. IEEE Transactions on Automatic Control, 1974.\par
A. Raue, C. Kreutz, T. Maiwald, et al. Structural and practical identifiability analysis of partially observed
dynamical models by exploiting the profile likelihood. Bioinformatics, 2009.\par
D. Prichard and J. Theiler. Generating surrogate data for time series with several simultaneously measured
variables. Physical Review Letters, 1994.\par
H. Turlier, D. A. Fedosov, B. Audoly, et al. Equilibrium physics breakdown reveals the active nature of red blood
cell flickering. Nature Physics, 2016.\par}
\note{Not presented; for questions.}
\end{frame}

% =================================================================================================== appendix
\appendix

\begin{frame}
\gtitleL{Appendix: fluctuations fix ratios, not stiffness}
\gbox{0.43}{1.6}{3.05}{\begin{gitems}
\item Scaling the moduli, viscosity and temperature together leaves the data unchanged
\item So the profile of $\kappa$ is flat when the temperature is free
\item The ratios $\kappa/\eta$, $\bar\sigma$ and $\bar\gamma$ are still identifiable: recovered within 4\% on
      simulated data
\end{gitems}}
\geqn{4.0}{0.45}{5.6}{(\kappa,\ \sigma,\ \gamma,\ \eta,\ T_\mathrm{eff})\ \longmapsto\ c\,(\kappa,\ \sigma,\ \gamma,\ \eta,\ T_\mathrm{eff})}
\gimg{3.75}{1.21}{5.85}{3.9}{gs_identifiability.png}
\note{With temperature fixed, kappa's profile likelihood has a minimum (orange); with temperature free it is
flat to one part in a billion (blue). sigma-bar = sigma R squared over kappa; gamma-bar = gamma R to the fourth
over kappa. Consequence: the published trap-induced stiffening is a variance ratio, so it is equally a fall in
forcing.}
\end{frame}

\begin{frame}
\gtitle{Appendix: the passive model is rejected, even after calibration}
\gimgw{1.42}{0.97}{7.17}{model_selection.png}
\gbox{0.43}{4.0}{9.1}{\begin{gitems}[14]
\item P loses in all 11 videos (left, $\Delta$AIC $+796$ to $+1479$); refitted passive simulations give 95\%
      quantiles of only 27--554, and on synthetic passive data P wins 3 of 4 times (right)
\item This is not proof of activity: a passive model with a broad spread of rates comes close, and there is no
      ATP-depleted control
\end{gitems}}
\note{With 25,000 periodogram ordinates any wrong model loses AIC by a lot, so I calibrated: three passive
generators simulated through the full observation model and refit exactly like the data. The four calibration
cells gave 1151 to 1479.}
\end{frame}

\begin{frame}
\gtitle{Appendix: the colored-forcing model also misfits every cell}
\gimgw{0.95}{0.97}{8.11}{misfit.png}
\gbox{0.43}{4.0}{9.1}{\begin{gitems}[14]
\item Model A fits better, but misfits every cell: data over model is above 1 below 2\,Hz and below 1 at
      3--30\,Hz (median over modes $q=2$--6)
\item So its fitted parameters are best-fit values of a wrong model, and the forcing needs a spread of time scales
\end{gitems}}
\note{A posterior predictive check gives chi-squared of 488 to 900 where 99 percent of simulations fall near 200
or below. A second colored term helps but does not fix it.}
\end{frame}

\begin{frame}
\gtitleL{Appendix: injection tests}
\gbox{0.43}{1.6}{3.05}{\begin{gitems}
\item 100\,nm kicks injected into real frames were detected in only 4 of 12 runs
\item Real cells: kurtosis 3.01--3.39, not the 4.1 reported for free cells
\item Tracking noise can turn an MSD exponent of 1 into the reported $\sim$1/4; corrected, it is 0.29--0.81 in
      7 cells
\end{gitems}}
\gimg{3.75}{0.88}{5.85}{3.9}{gs_markers.png}
\note{Caselli et al. report localized active kickers of about 100 nm. Narrow kicks (50 nm) vanish into the
Gaussian band in all 12 runs; at 100 nm they are detected in 4 of 12, so the reported kickers sit at the
detection limit. One video and a guessed kick rate, since none is reported.}
\end{frame}

\begin{frame}
\gtitleL{Appendix: a second-order effect}
\gbox{0.43}{1.6}{3.05}{\begin{gitems}
\item Particles on the fitted walls of one cell, scaled by a factor $s$
\item The effect follows $bs^2+cs^4$; a held-out run at $s=1$ agrees in 5 of 6 cases
\item Predicted effect at the real amplitude: $b=-0.14\pm0.06\%$ of $\Phi$ over 12 walls of this cell
\end{gitems}}
\gimg{3.75}{0.69}{5.85}{4.27}{gs_particles.png}
\note{The six cases are 3 walls times 2 target radii (0.3R and 0.5R). On the same walls the grid extrapolation
gives -0.08 +/- 0.04 percent; the paired difference is -0.06 +/- 0.03 percent. Both are of order 0.1 percent.}
\end{frame}

\end{document}
